\section{关键挑战：主路径如何相互放大}

国产推理引擎面临的难点并非单点性能不足，而是多个系统目标同时受限：既要持续吸收上游能力演进，又要适应本土硬件与基础软件栈的边界，还要在真实业务负载下维持稳定交付。这些挑战之所以顽固，是因为它们会相互放大——平台差异抬高抽象成本，抽象不稳又削弱复杂任务支持；复杂任务若采用侵入式补丁实现，又加速本地分叉累积；而如果评测仍主要停留在理想化 benchmark，团队就难以及时识别这些结构性问题。表\ref{tab:challenge-map} 将这种放大关系压缩为总览。

\begin{table}[t]
\centering
\small
\caption{国产推理引擎关键挑战比较}
\label{tab:challenge-map}
\begin{tabularx}{0.88\linewidth}{L{2.2cm}L{4.2cm}Y}
\toprule
挑战来源 & 主要冲击层 & 典型后果 \\
\midrule
硬件异构与后端碎片化 & 执行层、算子接口、通信路径 & 共享路径分叉、正确性边界不一致、维护成本持续上升 \\
复杂任务与多模态负载 & 调度、缓存、图模式与前处理链路 & 阶段争用加剧、时延波动放大、端到端不稳定 \\
压缩收益与系统成本错配 & 量化路径、KV 状态、调度与成本核算 & 低比特收益不稳定、单位 token 成本难形成闭环 \\
上游快速演进与本地分叉 & 平台抽象、插件接口与共享模块 & patch 累积、合并困难、对新模型和新特性的跟进滞后 \\
评测体系与真实业务脱节 & benchmark 设计、运维指标与资源评估 & 优化方向失真、离线结论与生产表现偏离 \\
\bottomrule
\end{tabularx}
\end{table}

\begin{table*}[t]
\centering
\footnotesize
\setlength{\tabcolsep}{4pt}
\caption{不同国产推理路线与关键挑战的映射关系}
\label{tab:route-challenge-mapping}
\begin{tabularx}{\textwidth}{L{2.1cm}L{2.8cm}L{2.5cm}L{2.8cm}L{2.8cm}Y}
\toprule
代表路线 & 版本矩阵/环境门槛 & 网络与分布式风险 & 缓存/调度压力 & 运维与回退焦点 & 更需要优先建设的系统能力 \\
\midrule
vLLM-Ascend 类开源后端适配 & 高。上游主干、插件分支、CANN 与 torch\_npu 对齐要求强\citep{vllmascend} & 中到高。PD 分离与多机连通性常在扩容时暴露 & 中。继承上游调度能力，但平台约束可能导致命中率下降 & 版本锁定与环境预检 & 持续回归基线与插件边界稳定性 \\
MindIE 类一体化平台 & 中。平台内控但对外兼容性需额外验证\citep{mindie} & 低到中。平台内一体治理 & 中。内部调度但缺乏外部可观测证据 & 模板敏感性与配置可控性 & 服务模板的可诊断性与生态兼容接口 \\
Chitu / xLLM 类自研路线 & 中。自主控制但生态距离远\citep{chitu_github,xllm_docs} & 取决于自研通信层成熟度 & 高。状态治理需自建 & 全栈自主但生态覆盖有限 & 模型覆盖速度与社区生态对接 \\
vLLM-HUST 类混合式路线 & 中到高。继承上游版本压力，同时维护本地插件层\citep{vllmhust} & 中。继承上游分布式能力，补充平台护栏 & 中。上游调度加本地门控 & 模块边界维护与 merge-safe 演进 & 持续经营模块边界防止双重分叉 \\
\bottomrule
\end{tabularx}
\end{table*}

表\ref{tab:route-challenge-mapping} 进一步将这种放大关系映射到具体路线上，说明不同路线并非面对同一组问题再做不同优化，而是会把“版本矩阵、网络控制面、缓存/调度、运维回退”这些挑战优先暴露在不同边界上。路线选择本身，也就在决定团队首先要把系统能力压在哪一层。

公开社区部署记录与工程经验反复印证了这一判断：CANN/torch\_npu 版本矩阵对齐困难（vLLM-Ascend 0.7.x 到 0.8.x 升级期间多次出现 HCCL 通信与算子兼容性问题）、PD 分离与双机组网的稳定性、MindIE 服务化配置可控性，以及 LMDeploy 与 vLLM 在国产 GPU 环境中的适配取舍——这些反复出现的摩擦点，最终都能回收到表\ref{tab:challenge-map} 列出的几类结构性挑战上。真实摩擦往往先出现在环境治理和平台工具链层面，而不是单点跑分上。下面沿前文三条主路径，逐一讨论这些挑战如何沿系统骨架扩散。

\begin{figure}[!htb]
\centering
\resizebox{0.80\linewidth}{!}{\definecolor{cBlue}{HTML}{E6F0FA}
\definecolor{cGreen}{HTML}{E9F5EE}
\definecolor{cOrange}{HTML}{FCEFE3}
\definecolor{cRed}{HTML}{F7E6E6}
\definecolor{cGray}{HTML}{F2F3F5}
\begin{tikzpicture}[
    node distance=11mm and 14mm,
    scale=1.10,
    transform shape,
    challenge/.style={draw=black!70, line width=0.8pt, rounded corners=2.2pt, minimum width=3.0cm, minimum height=10mm, align=center, font=\small, inner sep=4pt},
    flow/.style={-{Latex[length=2.5mm]}, line width=0.85pt, draw=black!70},
    amplify/.style={-{Latex[length=2.5mm]}, line width=0.9pt, draw=red!65, dashed},
    note/.style={font=\footnotesize, align=left, text=black!70}
]

% Center: unified base
\node[challenge, fill=cGray, minimum width=3.8cm, minimum height=10mm] (base) {\textbf{统一底座}\\上游演进与评测体系};

% Four challenges around
\node[challenge, fill=cBlue, above left=13mm and 4mm of base] (c1) {\textbf{平台异构扩散}\\执行路径持续分化};
\node[challenge, fill=cGreen, above right=13mm and 4mm of base] (c2) {\textbf{复杂任务冲击}\\调度/缓存/状态失控};
\node[challenge, fill=cOrange, below left=13mm and 4mm of base] (c3) {\textbf{压缩收益失调}\\离线收益难以兑现};
\node[challenge, fill=cRed, below right=13mm and 4mm of base] (c4) {\textbf{评测体系脱节}\\优化方向与实际背离};

% Flows from challenges to base (challenges feed into base)
\draw[flow] (c1) -- (base);
\draw[flow] (c2) -- (base);
\draw[flow] (c3) -- (base);
\draw[flow] (c4) -- (base);

% Amplification arrows between challenges
\draw[amplify, bend left=15] (c1) to node[above, font=\scriptsize, text=red!70] {平台差异抬高抽象成本} (c2);
\draw[amplify, bend left=15] (c2) to node[right, font=\scriptsize, text=red!70] {侵入补丁加速分叉} (c3);
\draw[amplify, bend left=15] (c3) to node[below, font=\scriptsize, text=red!70] {离线错觉掩盖缺陷} (c4);
\draw[amplify, bend left=15] (c4) to node[left, font=\scriptsize, text=red!70] {评测失真延误识别} (c1);

% Center labels
\node[note, text width=5cm, below=3mm of base] {四类挑战形成相互放大循环};
\node[note, text width=4cm, above=12mm of c1, xshift=28mm] {单一挑战不可孤立解决};
\end{tikzpicture}
}
\caption{关键挑战的相互放大关系：四类结构性问题通过统一底座形成闭环反馈}
\label{fig:challenge-amplification}
\end{figure}

\subsection{执行与通信主路径：平台异构的持续扩散}

第 3 节已将后端能力契约列为执行路径的核心问题，本节关注的是：为什么平台异构在国产推理引擎中不是一次性适配问题，而是沿执行路径持续扩散的系统风险。

平台异构的特殊难点在于硬件代际能力差异大且缺少统一适配层。vLLM-Ascend、vLLM-MLU 与 SGLang Ascend 的公开材料都体现出对上游主线版本、CANN SDK 和 torch\_npu 三方联合对齐的强依赖\citep{vllmascend,vllmmlu,sglang_ascend_npu}。不同硬件上的 feature parity 难以默认成立——graph mode、attention backend、量化、PD 分离和 speculative decoding 等能力即使在同一工具链中也可能存在支持粒度差异\citep{fastdeploy_docs,lmdeploy_supported_models}。公开性能数字不能直接横向比较，因为模型版本、输入/输出长度、并发度、精度和 runtime 版本的任一变化都可能改变结论。

\subsection{状态治理主路径：复杂任务场景的能力补齐}

当前真实负载已从单轮文本生成扩展到长上下文、多模态、结构化输出、工具调用和多 Agent 协同\citep{dong2024xgrammar,liu2025elasticmm,xie2025aimetropolis}。推理引擎的优化目标不再是单一路径 decode，而是跨阶段状态管理：结构化输出要求把约束执行压进采样环路，工具调用要求支持生成与外部执行之间的往返切换，多模态请求要求消化更长的前处理链和更剧烈的 shape 波动。对国产平台，难点不在接口名义上"可用"，而在这些能力会同时冲击调度、缓存和图模式。

复杂任务场景直接改写了"状态"的性质。文本 serving 时代系统重点管理请求队列、KV Cache 和 batch 形状；进入多模态和工具调用场景后，运行时还须管理视觉前缀、语法状态、工具回注上下文、外部执行等待与流式中断恢复等状态对象。它们既非纯粹的模型内部变量，也无法完全交给服务层处理，而是直接影响调度顺序、缓存复用和回退路径选择的运行时实体。

\begin{figure}[!htb]
\centering
\resizebox{0.85\linewidth}{!}{\begin{tikzpicture}[
    node distance=5mm and 4mm,
    stage/.style={draw, rounded corners=2pt, minimum width=22mm, minimum height=10mm, align=center, font=\small, inner sep=4pt},
    side/.style={draw, rounded corners=2pt, align=center, font=\scriptsize, inner sep=3pt},
    flow/.style={-{Latex[length=2.2mm]}, semithick}
]
\node[stage, fill=gray!10] (ingress) {请求进入\\文本/多模态/工具调用};
\node[stage, fill=blue!10, right=of ingress] (normalize) {请求整形\\长度 bucket / SLO};
\node[stage, fill=green!10, right=of normalize] (prefill) {前处理与 Prefill\\视觉/语音编码};
\node[stage, fill=green!10, right=of prefill] (decode) {Decode\\采样/约束执行};
\node[stage, fill=orange!10, right=of decode] (verify) {验证与回退\\grammar / speculative};
\node[stage, fill=gray!10, right=of verify] (finish) {响应返回\\状态释放与观测};

\node[side, fill=yellow!8, below=9mm of prefill] (cache) {共享状态\\KV Cache / prefix / 迁移元数据};
\node[side, fill=yellow!8, below=9mm of decode] (tool) {外部交互\\工具执行 / 回注 / 再进入};

\draw[flow] (ingress) -- (normalize);
\draw[flow] (normalize) -- (prefill);
\draw[flow] (prefill) -- (decode);
\draw[flow] (decode) -- (verify);
\draw[flow] (verify) -- (finish);
\draw[flow, dashed] (prefill) -- (cache);
\draw[flow, dashed] (decode) -- (cache);
\draw[flow, dashed] (verify) -- (cache);
\draw[flow, dashed] (decode) -- (tool);
\draw[flow, dashed] (tool.west) |- (normalize.south);

\node[font=\scriptsize, align=center, below=8mm of normalize] {复杂请求的关键难点不在单步执行，而在跨阶段状态传递、资源复用与回退一致性。};
\end{tikzpicture}}
\caption{复杂请求生命周期示意图}
\label{fig:complex-request-lifecycle}
\end{figure}

推测解码进一步抬高了对执行一致性的要求。SpecInfer、Medusa、EAGLE 与 Lookahead Decoding 分别代表辅助模型、多头预测、特征不确定性建模与并行前瞻路线\citep{specinfer2023,leviathan2023fast,cai2024medusa,li2024eagle,fu2024lookahead}。硬件侧经验表明收益高度条件化：在 MI300X 上 speculative decoding 主要受益于小 batch 的带宽受限 decode，batch 增大后收益明显收缩\citep{singh2025specdecmi300x}。这类机制对国产平台尤其敏感，因为本地硬件在图模式切换和动态 shape 支持上约束更强——一个"只是多一次验证"的机制，迁移到新后端后可能同时意味着图缓存失效和异常回退链变长。

\subsection{压缩与成本主路径：收益不稳定与闭环缺失}

压缩路径的挑战与执行路径和状态治理直接耦合。执行路径上的平台异构会限制低比特算子的可用边界，状态治理层的缓存组织则决定压缩策略能否在运行时稳定兑现——三者缺一，压缩收益就难以从离线实验转化为生产能力。

"低成本"不是部署阶段最后再追问的经济性指标，而是与执行效率和状态治理同级的系统目标。问题在于压缩策略的收益往往不线性兑现：同一组量化参数在不同后端、不同上下文长度和不同负载强度下，可能分别表现为吞吐提升、尾时延恶化或精度波动。QServe 说明只有把量化、kernel 组织和 serving runtime 一起设计，低比特收益才可能在真实服务路径中落地\citep{wang2024qserve}；KIVI 与 SnapKV 表明长上下文场景下的状态压缩如果不和恢复路径与质量门限一起考虑，容量节省就会转化为恢复时延或命中率损失\citep{liu2024kivi,li2024snapkv}。对国产平台，这种耦合还会被放大——许多后端对低比特算子、动态 shape 和异步搬运的支持边界不一致。

\subsection{统一底座：上游演进与评测体系}

执行、状态和压缩三条主路径上的挑战，最终都需要回到“统一底座”层面来解决。所谓统一底座，不仅指技术架构上的抽象与观测基础设施，更指系统能否在持续演进中维持工程秩序——这涉及上游演进策略与评测体系两个维度。

高性能推理框架仍在快速演化。国产后端如果通过侵入式修改共享路径获得短期收益，后续将面临高昂合并成本。本地分叉之所以危险，还因为它逐步改变团队的优化激励——若长期依赖复制上游模块并在本地修补，团队最容易积累的是"快速打补丁"而非"设计稳定扩展面"的经验。相反，若平台特化更多收敛在插件、注册表和受控配置门面中，团队就更有机会把适配劳动沉淀为可复用的工程机制。因此 merge-safe 不是"对上游友好"的附属要求，而是国产推理引擎能否持续演进的核心条件。

评测体系方面，当前不少优化仍主要围绕静态 benchmark 展开。但真实业务包含请求长度分布变化大、多租户并发、冷热模型混部和 SLA 约束严格等特征。若评测只覆盖理想化场景，就难以反映系统在生产环境中的优势与短板。未来评测需要从"单点性能测试"扩展为"端到端工作负载评测"，将 TTFT、稳定吞吐、缓存命中率、异常恢复时间和资源成本纳入统一分析框架。OpenCompass 与 FlagPerf 分别代表了模型能力评测和系统性能评测的生态努力\citep{opencompass,flagperf}。真正有价值的评测应当连接这两类视角——既看模型效果，也看支撑这些效果所付出的系统成本。评测体系需要分层：底层算子用 microbenchmark 验证，运行时调度用受控 serving workload 验证（重点是 TTFT、TPOT 与尾时延），而平台演进能力和运维代价则更适合通过持续集成基线、版本升级回归和故障演练来评估。
